\section{}

\paragraph{1150年}\eventanchor{SS-1150-REGNAL-YEAR}{SS-1150-REGNAL-YEAR}　南宋以宋高宗在位第24年作为诏令、赋役与外交文书的纪年框架。账本列出《宋史》卷24—47〈本纪〉；《建炎以来系年要录》本年条；《宋会要辑稿》相关门类作为来源线索；该材料可核对时间与制度称谓，但有关执行范围、群体经验、军数或后果，仍须与地方文书、物质遗存及相关政权记录互证。

\paragraph{1151年}\eventanchor{SS-1151-REGNAL-YEAR}{SS-1151-REGNAL-YEAR}　南宋以宋高宗在位第25年作为诏令、赋役与外交文书的纪年框架。账本列出《宋史》卷24—47〈本纪〉；《建炎以来系年要录》本年条；《宋会要辑稿》相关门类作为来源线索；该材料可核对时间与制度称谓，但有关执行范围、群体经验、军数或后果，仍须与地方文书、物质遗存及相关政权记录互证。

\paragraph{1152年}\eventanchor{SS-1152-REGNAL-YEAR}{SS-1152-REGNAL-YEAR}　南宋以宋高宗在位第26年作为诏令、赋役与外交文书的纪年框架。账本列出《宋史》卷24—47〈本纪〉；《建炎以来系年要录》本年条；《宋会要辑稿》相关门类作为来源线索；该材料可核对时间与制度称谓，但有关执行范围、群体经验、军数或后果，仍须与地方文书、物质遗存及相关政权记录互证。

\paragraph{1153年}\eventanchor{SS-1153-REGNAL-YEAR}{SS-1153-REGNAL-YEAR}　南宋以宋高宗在位第27年作为诏令、赋役与外交文书的纪年框架。账本列出《宋史》卷24—47〈本纪〉；《建炎以来系年要录》本年条；《宋会要辑稿》相关门类作为来源线索；该材料可核对时间与制度称谓，但有关执行范围、群体经验、军数或后果，仍须与地方文书、物质遗存及相关政权记录互证。

\paragraph{1154年}\eventanchor{SS-1154-REGNAL-YEAR}{SS-1154-REGNAL-YEAR}　南宋以宋高宗在位第28年作为诏令、赋役与外交文书的纪年框架。账本列出《宋史》卷24—47〈本纪〉；《建炎以来系年要录》本年条；《宋会要辑稿》相关门类作为来源线索；该材料可核对时间与制度称谓，但有关执行范围、群体经验、军数或后果，仍须与地方文书、物质遗存及相关政权记录互证。

\paragraph{1155年}\eventanchor{SS-1155-REGNAL-YEAR}{SS-1155-REGNAL-YEAR}　南宋以宋高宗在位第29年作为诏令、赋役与外交文书的纪年框架。账本列出《宋史》卷24—47〈本纪〉；《建炎以来系年要录》本年条；《宋会要辑稿》相关门类作为来源线索；该材料可核对时间与制度称谓，但有关执行范围、群体经验、军数或后果，仍须与地方文书、物质遗存及相关政权记录互证。

\paragraph{1156年}\eventanchor{SS-1156-REGNAL-YEAR}{SS-1156-REGNAL-YEAR}　南宋以宋高宗在位第30年作为诏令、赋役与外交文书的纪年框架。账本列出《宋史》卷24—47〈本纪〉；《建炎以来系年要录》本年条；《宋会要辑稿》相关门类作为来源线索；该材料可核对时间与制度称谓，但有关执行范围、群体经验、军数或后果，仍须与地方文书、物质遗存及相关政权记录互证。

\paragraph{1157年}\eventanchor{SS-1157-REGNAL-YEAR}{SS-1157-REGNAL-YEAR}　南宋以宋高宗在位第31年作为诏令、赋役与外交文书的纪年框架。账本列出《宋史》卷24—47〈本纪〉；《建炎以来系年要录》本年条；《宋会要辑稿》相关门类作为来源线索；该材料可核对时间与制度称谓，但有关执行范围、群体经验、军数或后果，仍须与地方文书、物质遗存及相关政权记录互证。

\paragraph{1158年}\eventanchor{SS-1158-REGNAL-YEAR}{SS-1158-REGNAL-YEAR}　南宋以宋高宗在位第32年作为诏令、赋役与外交文书的纪年框架。账本列出《宋史》卷24—47〈本纪〉；《建炎以来系年要录》本年条；《宋会要辑稿》相关门类作为来源线索；该材料可核对时间与制度称谓，但有关执行范围、群体经验、军数或后果，仍须与地方文书、物质遗存及相关政权记录互证。

\paragraph{1159年}\eventanchor{SS-1159-REGNAL-YEAR}{SS-1159-REGNAL-YEAR}　南宋以宋高宗在位第33年作为诏令、赋役与外交文书的纪年框架。账本列出《宋史》卷24—47〈本纪〉；《建炎以来系年要录》本年条；《宋会要辑稿》相关门类作为来源线索；该材料可核对时间与制度称谓，但有关执行范围、群体经验、军数或后果，仍须与地方文书、物质遗存及相关政权记录互证。

\paragraph{1160年}\eventanchor{SS-1160-REGNAL-YEAR}{SS-1160-REGNAL-YEAR}　南宋以宋高宗在位第34年作为诏令、赋役与外交文书的纪年框架。账本列出《宋史》卷24—47〈本纪〉；《建炎以来系年要录》本年条；《宋会要辑稿》相关门类作为来源线索；该材料可核对时间与制度称谓，但有关执行范围、群体经验、军数或后果，仍须与地方文书、物质遗存及相关政权记录互证。

\paragraph{1161年}\eventanchor{SS-1161-EVENT-123}{SS-1161-EVENT-123}　南宋在采石等地击退金军渡江攻势。账本列出《宋史》卷24—47〈本纪〉；《建炎以来系年要录》本年条；《宋会要辑稿》相关门类作为来源线索；该材料可核对时间与制度称谓，但有关执行范围、群体经验、军数或后果，仍须与地方文书、物质遗存及相关政权记录互证。

\paragraph{1161年}\eventanchor{SS-1161-REGNAL-YEAR}{SS-1161-REGNAL-YEAR}　南宋以宋高宗在位第35年作为诏令、赋役与外交文书的纪年框架。账本列出《宋史》卷24—47〈本纪〉；《建炎以来系年要录》本年条；《宋会要辑稿》相关门类作为来源线索；该材料可核对时间与制度称谓，但有关执行范围、群体经验、军数或后果，仍须与地方文书、物质遗存及相关政权记录互证。

\paragraph{1162年}\eventanchor{SS-1162-EVENT-124}{SS-1162-EVENT-124}　高宗禅位，孝宗即位。账本列出《宋史》卷24—47〈本纪〉；《建炎以来系年要录》本年条；《宋会要辑稿》相关门类作为来源线索；该材料可核对时间与制度称谓，但有关执行范围、群体经验、军数或后果，仍须与地方文书、物质遗存及相关政权记录互证。

\paragraph{1162年}\eventanchor{SS-1162-REGNAL-YEAR}{SS-1162-REGNAL-YEAR}　南宋以宋高宗在位第36年作为诏令、赋役与外交文书的纪年框架。账本列出《宋史》卷24—47〈本纪〉；《建炎以来系年要录》本年条；《宋会要辑稿》相关门类作为来源线索；该材料可核对时间与制度称谓，但有关执行范围、群体经验、军数或后果，仍须与地方文书、物质遗存及相关政权记录互证。

\paragraph{1163年}\eventanchor{SS-1163-EVENT-125}{SS-1163-EVENT-125}　孝宗支持北伐，宋军在符离等地受挫。账本列出《宋史》卷24—47〈本纪〉；《建炎以来系年要录》本年条；《宋会要辑稿》相关门类作为来源线索；该材料可核对时间与制度称谓，但有关执行范围、群体经验、军数或后果，仍须与地方文书、物质遗存及相关政权记录互证。

\paragraph{1163年}\eventanchor{SS-1163-REGNAL-YEAR}{SS-1163-REGNAL-YEAR}　南宋以宋孝宗在位第1年作为诏令、赋役与外交文书的纪年框架。账本列出《宋史》卷24—47〈本纪〉；《建炎以来系年要录》本年条；《宋会要辑稿》相关门类作为来源线索；该材料可核对时间与制度称谓，但有关执行范围、群体经验、军数或后果，仍须与地方文书、物质遗存及相关政权记录互证。

\paragraph{1164年}\eventanchor{SS-1164-EVENT-126}{SS-1164-EVENT-126}　宋金订立隆兴和议。账本列出《宋史》卷24—47〈本纪〉；《建炎以来系年要录》本年条；《宋会要辑稿》相关门类作为来源线索；该材料可核对时间与制度称谓，但有关执行范围、群体经验、军数或后果，仍须与地方文书、物质遗存及相关政权记录互证。

\paragraph{1164年}\eventanchor{SS-1164-REGNAL-YEAR}{SS-1164-REGNAL-YEAR}　南宋以宋孝宗在位第2年作为诏令、赋役与外交文书的纪年框架。账本列出《宋史》卷24—47〈本纪〉；《建炎以来系年要录》本年条；《宋会要辑稿》相关门类作为来源线索；该材料可核对时间与制度称谓，但有关执行范围、群体经验、军数或后果，仍须与地方文书、物质遗存及相关政权记录互证。

\paragraph{1165年}\eventanchor{SS-1165-REGNAL-YEAR}{SS-1165-REGNAL-YEAR}　南宋以宋孝宗在位第3年作为诏令、赋役与外交文书的纪年框架。账本列出《宋史》卷24—47〈本纪〉；《建炎以来系年要录》本年条；《宋会要辑稿》相关门类作为来源线索；该材料可核对时间与制度称谓，但有关执行范围、群体经验、军数或后果，仍须与地方文书、物质遗存及相关政权记录互证。

\paragraph{1166年}\eventanchor{SS-1166-REGNAL-YEAR}{SS-1166-REGNAL-YEAR}　南宋以宋孝宗在位第4年作为诏令、赋役与外交文书的纪年框架。账本列出《宋史》卷24—47〈本纪〉；《建炎以来系年要录》本年条；《宋会要辑稿》相关门类作为来源线索；该材料可核对时间与制度称谓，但有关执行范围、群体经验、军数或后果，仍须与地方文书、物质遗存及相关政权记录互证。

\paragraph{1167年}\eventanchor{SS-1167-REGNAL-YEAR}{SS-1167-REGNAL-YEAR}　南宋以宋孝宗在位第5年作为诏令、赋役与外交文书的纪年框架。账本列出《宋史》卷24—47〈本纪〉；《建炎以来系年要录》本年条；《宋会要辑稿》相关门类作为来源线索；该材料可核对时间与制度称谓，但有关执行范围、群体经验、军数或后果，仍须与地方文书、物质遗存及相关政权记录互证。

\paragraph{1168年}\eventanchor{SS-1168-REGNAL-YEAR}{SS-1168-REGNAL-YEAR}　南宋以宋孝宗在位第6年作为诏令、赋役与外交文书的纪年框架。账本列出《宋史》卷24—47〈本纪〉；《建炎以来系年要录》本年条；《宋会要辑稿》相关门类作为来源线索；该材料可核对时间与制度称谓，但有关执行范围、群体经验、军数或后果，仍须与地方文书、物质遗存及相关政权记录互证。

\paragraph{1169年}\eventanchor{SS-1169-REGNAL-YEAR}{SS-1169-REGNAL-YEAR}　南宋以宋孝宗在位第7年作为诏令、赋役与外交文书的纪年框架。账本列出《宋史》卷24—47〈本纪〉；《建炎以来系年要录》本年条；《宋会要辑稿》相关门类作为来源线索；该材料可核对时间与制度称谓，但有关执行范围、群体经验、军数或后果，仍须与地方文书、物质遗存及相关政权记录互证。

\paragraph{1170年}\eventanchor{SS-1170-REGNAL-YEAR}{SS-1170-REGNAL-YEAR}　南宋以宋孝宗在位第8年作为诏令、赋役与外交文书的纪年框架。账本列出《宋史》卷24—47〈本纪〉；《建炎以来系年要录》本年条；《宋会要辑稿》相关门类作为来源线索；该材料可核对时间与制度称谓，但有关执行范围、群体经验、军数或后果，仍须与地方文书、物质遗存及相关政权记录互证。

\paragraph{1171年}\eventanchor{SS-1171-REGNAL-YEAR}{SS-1171-REGNAL-YEAR}　南宋以宋孝宗在位第9年作为诏令、赋役与外交文书的纪年框架。账本列出《宋史》卷24—47〈本纪〉；《建炎以来系年要录》本年条；《宋会要辑稿》相关门类作为来源线索；该材料可核对时间与制度称谓，但有关执行范围、群体经验、军数或后果，仍须与地方文书、物质遗存及相关政权记录互证。

\paragraph{1172年}\eventanchor{SS-1172-REGNAL-YEAR}{SS-1172-REGNAL-YEAR}　南宋以宋孝宗在位第10年作为诏令、赋役与外交文书的纪年框架。账本列出《宋史》卷24—47〈本纪〉；《建炎以来系年要录》本年条；《宋会要辑稿》相关门类作为来源线索；该材料可核对时间与制度称谓，但有关执行范围、群体经验、军数或后果，仍须与地方文书、物质遗存及相关政权记录互证。

\paragraph{1173年}\eventanchor{SS-1173-EVENT-127}{SS-1173-EVENT-127}　朱熹等在书院、讲学与地方事务中活动。账本列出《宋史》卷24—47〈本纪〉；《建炎以来系年要录》本年条；《宋会要辑稿》相关门类作为来源线索；该材料可核对时间与制度称谓，但有关执行范围、群体经验、军数或后果，仍须与地方文书、物质遗存及相关政权记录互证。

\paragraph{1173年}\eventanchor{SS-1173-REGNAL-YEAR}{SS-1173-REGNAL-YEAR}　南宋以宋孝宗在位第11年作为诏令、赋役与外交文书的纪年框架。账本列出《宋史》卷24—47〈本纪〉；《建炎以来系年要录》本年条；《宋会要辑稿》相关门类作为来源线索；该材料可核对时间与制度称谓，但有关执行范围、群体经验、军数或后果，仍须与地方文书、物质遗存及相关政权记录互证。

\paragraph{1174年}\eventanchor{SS-1174-REGNAL-YEAR}{SS-1174-REGNAL-YEAR}　南宋以宋孝宗在位第12年作为诏令、赋役与外交文书的纪年框架。账本列出《宋史》卷24—47〈本纪〉；《建炎以来系年要录》本年条；《宋会要辑稿》相关门类作为来源线索；该材料可核对时间与制度称谓，但有关执行范围、群体经验、军数或后果，仍须与地方文书、物质遗存及相关政权记录互证。

\paragraph{1175年}\eventanchor{SS-1175-REGNAL-YEAR}{SS-1175-REGNAL-YEAR}　南宋以宋孝宗在位第13年作为诏令、赋役与外交文书的纪年框架。账本列出《宋史》卷24—47〈本纪〉；《建炎以来系年要录》本年条；《宋会要辑稿》相关门类作为来源线索；该材料可核对时间与制度称谓，但有关执行范围、群体经验、军数或后果，仍须与地方文书、物质遗存及相关政权记录互证。

\paragraph{1176年}\eventanchor{SS-1176-REGNAL-YEAR}{SS-1176-REGNAL-YEAR}　南宋以宋孝宗在位第14年作为诏令、赋役与外交文书的纪年框架。账本列出《宋史》卷24—47〈本纪〉；《建炎以来系年要录》本年条；《宋会要辑稿》相关门类作为来源线索；该材料可核对时间与制度称谓，但有关执行范围、群体经验、军数或后果，仍须与地方文书、物质遗存及相关政权记录互证。

\paragraph{1177年}\eventanchor{SS-1177-REGNAL-YEAR}{SS-1177-REGNAL-YEAR}　南宋以宋孝宗在位第15年作为诏令、赋役与外交文书的纪年框架。账本列出《宋史》卷24—47〈本纪〉；《建炎以来系年要录》本年条；《宋会要辑稿》相关门类作为来源线索；该材料可核对时间与制度称谓，但有关执行范围、群体经验、军数或后果，仍须与地方文书、物质遗存及相关政权记录互证。

\paragraph{1178年}\eventanchor{SS-1178-REGNAL-YEAR}{SS-1178-REGNAL-YEAR}　南宋以宋孝宗在位第16年作为诏令、赋役与外交文书的纪年框架。账本列出《宋史》卷24—47〈本纪〉；《建炎以来系年要录》本年条；《宋会要辑稿》相关门类作为来源线索；该材料可核对时间与制度称谓，但有关执行范围、群体经验、军数或后果，仍须与地方文书、物质遗存及相关政权记录互证。

\paragraph{1179年}\eventanchor{SS-1179-REGNAL-YEAR}{SS-1179-REGNAL-YEAR}　南宋以宋孝宗在位第17年作为诏令、赋役与外交文书的纪年框架。账本列出《宋史》卷24—47〈本纪〉；《建炎以来系年要录》本年条；《宋会要辑稿》相关门类作为来源线索；该材料可核对时间与制度称谓，但有关执行范围、群体经验、军数或后果，仍须与地方文书、物质遗存及相关政权记录互证。

\paragraph{1180年}\eventanchor{SS-1180-REGNAL-YEAR}{SS-1180-REGNAL-YEAR}　南宋以宋孝宗在位第18年作为诏令、赋役与外交文书的纪年框架。账本列出《宋史》卷24—47〈本纪〉；《建炎以来系年要录》本年条；《宋会要辑稿》相关门类作为来源线索；该材料可核对时间与制度称谓，但有关执行范围、群体经验、军数或后果，仍须与地方文书、物质遗存及相关政权记录互证。

\paragraph{1181年}\eventanchor{SS-1181-REGNAL-YEAR}{SS-1181-REGNAL-YEAR}　南宋以宋孝宗在位第19年作为诏令、赋役与外交文书的纪年框架。账本列出《宋史》卷24—47〈本纪〉；《建炎以来系年要录》本年条；《宋会要辑稿》相关门类作为来源线索；该材料可核对时间与制度称谓，但有关执行范围、群体经验、军数或后果，仍须与地方文书、物质遗存及相关政权记录互证。

\paragraph{1182年}\eventanchor{SS-1182-REGNAL-YEAR}{SS-1182-REGNAL-YEAR}　南宋以宋孝宗在位第20年作为诏令、赋役与外交文书的纪年框架。账本列出《宋史》卷24—47〈本纪〉；《建炎以来系年要录》本年条；《宋会要辑稿》相关门类作为来源线索；该材料可核对时间与制度称谓，但有关执行范围、群体经验、军数或后果，仍须与地方文书、物质遗存及相关政权记录互证。

\paragraph{1183年}\eventanchor{SS-1183-REGNAL-YEAR}{SS-1183-REGNAL-YEAR}　南宋以宋孝宗在位第21年作为诏令、赋役与外交文书的纪年框架。账本列出《宋史》卷24—47〈本纪〉；《建炎以来系年要录》本年条；《宋会要辑稿》相关门类作为来源线索；该材料可核对时间与制度称谓，但有关执行范围、群体经验、军数或后果，仍须与地方文书、物质遗存及相关政权记录互证。

\paragraph{1184年}\eventanchor{SS-1184-REGNAL-YEAR}{SS-1184-REGNAL-YEAR}　南宋以宋孝宗在位第22年作为诏令、赋役与外交文书的纪年框架。账本列出《宋史》卷24—47〈本纪〉；《建炎以来系年要录》本年条；《宋会要辑稿》相关门类作为来源线索；该材料可核对时间与制度称谓，但有关执行范围、群体经验、军数或后果，仍须与地方文书、物质遗存及相关政权记录互证。

\paragraph{1185年}\eventanchor{SS-1185-REGNAL-YEAR}{SS-1185-REGNAL-YEAR}　南宋以宋孝宗在位第23年作为诏令、赋役与外交文书的纪年框架。账本列出《宋史》卷24—47〈本纪〉；《建炎以来系年要录》本年条；《宋会要辑稿》相关门类作为来源线索；该材料可核对时间与制度称谓，但有关执行范围、群体经验、军数或后果，仍须与地方文书、物质遗存及相关政权记录互证。

\paragraph{1186年}\eventanchor{SS-1186-REGNAL-YEAR}{SS-1186-REGNAL-YEAR}　南宋以宋孝宗在位第24年作为诏令、赋役与外交文书的纪年框架。账本列出《宋史》卷24—47〈本纪〉；《建炎以来系年要录》本年条；《宋会要辑稿》相关门类作为来源线索；该材料可核对时间与制度称谓，但有关执行范围、群体经验、军数或后果，仍须与地方文书、物质遗存及相关政权记录互证。

\paragraph{1187年}\eventanchor{SS-1187-REGNAL-YEAR}{SS-1187-REGNAL-YEAR}　南宋以宋孝宗在位第25年作为诏令、赋役与外交文书的纪年框架。账本列出《宋史》卷24—47〈本纪〉；《建炎以来系年要录》本年条；《宋会要辑稿》相关门类作为来源线索；该材料可核对时间与制度称谓，但有关执行范围、群体经验、军数或后果，仍须与地方文书、物质遗存及相关政权记录互证。

\paragraph{1188年}\eventanchor{SS-1188-REGNAL-YEAR}{SS-1188-REGNAL-YEAR}　南宋以宋孝宗在位第26年作为诏令、赋役与外交文书的纪年框架。账本列出《宋史》卷24—47〈本纪〉；《建炎以来系年要录》本年条；《宋会要辑稿》相关门类作为来源线索；该材料可核对时间与制度称谓，但有关执行范围、群体经验、军数或后果，仍须与地方文书、物质遗存及相关政权记录互证。

\paragraph{1189年}\eventanchor{SS-1189-EVENT-128}{SS-1189-EVENT-128}　孝宗去世后广宗即位。账本列出《宋史》卷24—47〈本纪〉；《建炎以来系年要录》本年条；《宋会要辑稿》相关门类作为来源线索；该材料可核对时间与制度称谓，但有关执行范围、群体经验、军数或后果，仍须与地方文书、物质遗存及相关政权记录互证。

\paragraph{1189年}\eventanchor{SS-1189-REGNAL-YEAR}{SS-1189-REGNAL-YEAR}　南宋以宋孝宗在位第27年作为诏令、赋役与外交文书的纪年框架。账本列出《宋史》卷24—47〈本纪〉；《建炎以来系年要录》本年条；《宋会要辑稿》相关门类作为来源线索；该材料可核对时间与制度称谓，但有关执行范围、群体经验、军数或后果，仍须与地方文书、物质遗存及相关政权记录互证。

\paragraph{1190年}\eventanchor{SS-1190-REGNAL-YEAR}{SS-1190-REGNAL-YEAR}　南宋以宋光宗在位第1年作为诏令、赋役与外交文书的纪年框架。账本列出《宋史》卷24—47〈本纪〉；《建炎以来系年要录》本年条；《宋会要辑稿》相关门类作为来源线索；该材料可核对时间与制度称谓，但有关执行范围、群体经验、军数或后果，仍须与地方文书、物质遗存及相关政权记录互证。

\paragraph{1191年}\eventanchor{SS-1191-REGNAL-YEAR}{SS-1191-REGNAL-YEAR}　南宋以宋光宗在位第2年作为诏令、赋役与外交文书的纪年框架。账本列出《宋史》卷24—47〈本纪〉；《建炎以来系年要录》本年条；《宋会要辑稿》相关门类作为来源线索；该材料可核对时间与制度称谓，但有关执行范围、群体经验、军数或后果，仍须与地方文书、物质遗存及相关政权记录互证。

\paragraph{1192年}\eventanchor{SS-1192-REGNAL-YEAR}{SS-1192-REGNAL-YEAR}　南宋以宋光宗在位第3年作为诏令、赋役与外交文书的纪年框架。账本列出《宋史》卷24—47〈本纪〉；《建炎以来系年要录》本年条；《宋会要辑稿》相关门类作为来源线索；该材料可核对时间与制度称谓，但有关执行范围、群体经验、军数或后果，仍须与地方文书、物质遗存及相关政权记录互证。

\paragraph{1193年}\eventanchor{SS-1193-REGNAL-YEAR}{SS-1193-REGNAL-YEAR}　南宋以宋光宗在位第4年作为诏令、赋役与外交文书的纪年框架。账本列出《宋史》卷24—47〈本纪〉；《建炎以来系年要录》本年条；《宋会要辑稿》相关门类作为来源线索；该材料可核对时间与制度称谓，但有关执行范围、群体经验、军数或后果，仍须与地方文书、物质遗存及相关政权记录互证。

